\documentclass{article}
\usepackage[slovene]{babel}
\usepackage{biblatex}
\addbibresource{literatura.bib}
% [
% sortcites=true, 
% %sorting=ynt,
% style=alphabetic
% url=false,
% doi=false,
% ]


\usepackage{hyperref}


\title{\LaTeX{} 3. - Sklicevanje in Literatura}
\author{Antonio Montero}

\begin{document}
\maketitle

\section*{Sklicevanje}
\label{sec:sklicevanje}


% S uporabo ukaza \texttt{\textbackslash label\{sec:sklicevanje\}} smo označili ta razdelek, nanj pa se lahko sklicujemo z ukazom \texttt{\textbackslash ref\{sec:sklicevanje\}}. 

% Na primer, v razdelku \nameref{sec:sklicevanje} smo predstavili sklicevanje na različne elemente v dokumentu.
% Na primer, v razdelku \cref{sec:sklicevanje} smo predstavili sklicevanje na različne elemente v dokumentu.

\section{Literatura}
% Prva knjiga, ki sem jo prebral, je bila \cite{hp1}. 
% % ...
% \begin{thebibliography}{9}
%   \bibitem{hp1} J. K. Rowling, \textit{Harry Potter in kamen modrosti}, Založba EPTA, 1997.
%   \bibitem{hp2} J. K. Rowling, \textit{Harry Potter in dvorana skrivnosti}, Založba EPTA, 1998.
% \end{thebibliography}


\vspace{1cm}

Avtor knjige \cite{dm2018} je moj sosed v pisarni.


\ldots{}, ki ga je prikazala že ga. Curie
v \cite{curie}
\vspace{1cm}

\ldots{} Pred kratkim sem se učil uporabljati Bib\LaTeX{} iz \cite{lshort}. Zdi se mi zelo uporabno. \ldots
\vspace{1cm}

\ldots{} to je mogoče preprosto razložiti z
dejstvom, da je bil Einstein časovni
popotnik \cite{dream} \ldots
\vspace{1cm}

Vedno jih moramo navajati tako, kot so navedeni v bibliografiji \cite{dream, lshort, hp1, curie}.






\clearpage
\nocite{*}

% \printbibliography[title={Viri in literatura}]

% \printbibliography[type=article,title={Samo članki}]

%  \printbibliography[type=book,title={Vse knjige}]

%   \printbibliography[keyword={physics},title={Reference, povezane s Fiziko}]
%  \printbibliography[notkeyword={latex},title={Vse, razen tistih, povezane z LaTeX-om}]
 
 
%  \printbibliography[keyword={matematika},title={Reference, povezane s Matematiko}]

 
%  \defbibfilter{custom}{keyword=physics or keyword=latex}     
%  \printbibliography[filter=custom,title={Reference, povezane s fiziko ali \LaTeX-om}]


% \printbibliography[title={Literatura, privzeto}]
% \printbibliography[title={Literatura, v kazalu}, heading=bibintoc]
% \printbibliography[title={Literatura, kot podrazdelek -- z zvezdico --}, heading=subbibliography]
% \printbibliography[title={Literatura kot razdelek, s številko}, heading=bibnumbered]
% \printbibliography[title={Literatura kot podrazdelek, brez št.-a, v kazalu}, heading=subbibintoc]
% \printbibliography[title={Literatura kot podrazdelek, s št.-o, v kazalu}, heading=subbibnumbered]
% \printbibliography[title={Čeprav napišem naslov, se nič ne natisne}, heading=none]
\end{document}